\documentclass[11pt,a4paper]{article}

\usepackage[margin=1.05in]{geometry}
\usepackage{amsmath,amssymb,amsthm}
\usepackage{mathtools}
\usepackage{booktabs}
\usepackage{listings}
\usepackage{xcolor}
\usepackage[colorlinks=true,linkcolor=black,citecolor=black,urlcolor=blue]{hyperref}
\usepackage{enumitem}
\usepackage[T1]{fontenc}
\usepackage[protrusion=true,expansion=false]{microtype}

%%%%%%%%%%%%%%%%%%%%%%%%%%%%%%%%%%%%%%%%%%%%%%%%%%%%%%%%%%%%%%%%%%%%%%%%%%%%%%%
% Theorem environments
%%%%%%%%%%%%%%%%%%%%%%%%%%%%%%%%%%%%%%%%%%%%%%%%%%%%%%%%%%%%%%%%%%%%%%%%%%%%%%%
\theoremstyle{plain}
\newtheorem{theorem}{Theorem}[section]
\newtheorem{proposition}[theorem]{Proposition}
\newtheorem{lemma}[theorem]{Lemma}
\newtheorem{corollary}[theorem]{Corollary}
\newtheorem{conjecture}[theorem]{Conjecture}

\theoremstyle{definition}
\newtheorem{definition}[theorem]{Definition}
\newtheorem{example}[theorem]{Example}
\newtheorem{principle}[theorem]{Principle}
\newtheorem{obligation}[theorem]{Obligation}

\theoremstyle{remark}
\newtheorem{remark}[theorem]{Remark}

%%%%%%%%%%%%%%%%%%%%%%%%%%%%%%%%%%%%%%%%%%%%%%%%%%%%%%%%%%%%%%%%%%%%%%%%%%%%%%%
% Macros.  NOTATION RULES OBSERVED THROUGHOUT:
%   * the rho calculus is never written with the Greek letter rho
%   * quoting is @P, dropping is *x; corner quotes are pre-2016 notation
%%%%%%%%%%%%%%%%%%%%%%%%%%%%%%%%%%%%%%%%%%%%%%%%%%%%%%%%%%%%%%%%%%%%%%%%%%%%%%%
\newcommand{\rhoc}{\textsf{rho}}
\newcommand{\grho}{\textsf{rho}_{\V}}
\newcommand{\sem}[1]{[\![#1]\!]}
\newcommand{\V}{\mathbb{V}}
\newcommand{\Bool}{\mathbb{B}}
\newcommand{\Rnn}{\mathbb{R}_{\geq 0}}
\newcommand{\Cx}{\mathbb{C}}
\newcommand{\Nat}{\mathbb{N}}
\newcommand{\Proc}{\mathcal{P}}
\newcommand{\Names}{\mathcal{N}}
\newcommand{\NF}{\mathrm{NF}}
\newcommand{\Reach}{\mathrm{Reach}}
\newcommand{\supp}{\mathrm{supp}}
\newcommand{\Per}{\mathrm{Per}}
\newcommand{\Det}{\mathrm{Det}}
\newcommand{\scong}{\equiv}
\newcommand{\tensorv}{\otimes}
\newcommand{\sumv}{\oplus}
\newcommand{\Cand}{\mathrm{Cand}}
\newcommand{\Out}{\mathrm{out}}
\newcommand{\crisp}[1]{\lceil #1 \rceil}
\newcommand{\pow}{\mathrm{pow}}
\newcommand{\str}{\mathrm{s}}
\newcommand{\conf}{\mathrm{c}}

\lstdefinestyle{rholang}{
  basicstyle=\ttfamily\small,
  keywordstyle=\bfseries,
  morekeywords={for,new,in,where,contract,match,if,else},
  columns=fullflexible,
  keepspaces=true,
  showstringspaces=false,
  breaklines=true,
  frame=single,
  framesep=5pt,
  rulecolor=\color{black!35},
  xleftmargin=6pt,
  aboveskip=9pt,
  belowskip=9pt
}
\lstset{style=rholang}

%%%%%%%%%%%%%%%%%%%%%%%%%%%%%%%%%%%%%%%%%%%%%%%%%%%%%%%%%%%%%%%%%%%%%%%%%%%%%%%

\title{\bfseries Graded Where-Clauses\\[4pt]
\large Non-crisp truth values as the syntax of non-determinism resolution\\
in the rho calculus, with quantum and probabilistic-logic instances}

\author{L.\,G.\ Meredith\\ \small F1R3FLY.io}

\date{F1R3FLY.io working note --- August 2026}

\begin{document}
\maketitle

\begin{abstract}
The rho calculus is silent about how the non-determinism of a race is resolved.
An implementation must supply that mechanism, and different mechanisms give
different executions of the same program: RSpace supplies the physical races of
cores on a CPU, and one may equally supply the Born rule. We take this seriously
as a design principle and ask what syntax a programmer needs in order to
\emph{arrange} a resolution mechanism rather than merely inherit one.

We argue, first, that no repair of full abstraction over rho contexts can expose
the additional structure a quantum resolver supplies, because no rho context says
anything about how non-determinism is resolved; the mismatch is not in the class
of contexts but in the type of observation. We then show that unmodified rho,
read as a sum over its possible futures, already interferes --- but only on
structurally indistinguishable racers, and only constructively, because the
message left behind by a race is a which-path record.

The design that follows is a single change: the \texttt{where} clause of the
for-comprehension is valued in a commutative semiring $\V$ rather than in the
Booleans, and is evaluated \emph{per candidate match} rather than per
communication. Everything else is unchanged. We show that (i) weight maps in the
sense of the weighted-GSLT programme are exactly the linear normal forms of
graded clauses, so the partition discipline and the update-function machinery
both disappear; (ii) the join operator \texttt{\&} is already a recombiner, so
the which-path erasure that interference requires needs no new term former; and
(iii) the graded clause on an $m$-fold join over $m$ data \emph{is} a scattering
matrix, with the amplitude of the symmetric outcome given by its permanent.
Boolean $\V$ recovers rholang 1.4 exactly; $\Rnn$ gives stochastic execution;
$\Cx$ gives interference, including cancellation of an operationally reachable
outcome.

As a worked example we build Shor's algorithm in the fragment, with every number
computed by a simulator that enumerates candidates rather than multiplying gate
matrices. Universality is easy and modular exponentiation is ordinary rholang,
but the exercise turns up a hazard rather than a triumph: with clauses
unconstrained, renormalising at observation is post-selection, and the fragment
computes $\mathsf{PP}$. We therefore state the isometry condition that contains
it, and observe that the fragment on which it is decidable is exactly the
fragment a compiler to circuits can accept.

Finally we observe that the same slot accepts truth values in the sense of
Goertzel's Probabilistic Logic Networks. Valuing where-clauses by PLN inference
makes MeTTaIL's execution engine and its reasoning engine the same loop: what
fires next is decided by what the system currently believes. We give the
resolution algebra this requires, and are explicit about which part of PLN's
truth-value arithmetic supplies values and which part would have to be weakened
to supply the algebra. A second worked example --- a mortal forager whose
hypothesis is a PLN truth value in a \texttt{where} guard --- shows what the
grading buys: a Boolean guard has no coordinate in which to record
disappointment and cannot learn to stop, exploration and exploitation cease to
be scheduled, and PLN's personality parameter, elsewhere set by taste, acquires
an interior optimum fixed by the risk of the environment.
\end{abstract}

\tableofcontents

%%%%%%%%%%%%%%%%%%%%%%%%%%%%%%%%%%%%%%%%%%%%%%%%%%%%%%%%%%%%%%%%%%%%%%%%%%%%%%%
\section{Introduction}\label{sec:intro}

\subsection{Where the choice is}

A rho program does not determine its own execution. Write
\[
  x!(Q_1) \mid \textsf{for}(y \Leftarrow x)P \mid x!(Q_2)
\]
and the operational semantics permits two reductions. It does not say which one
happens, and it does not say with what tendency. This is not an oversight; it is
the defining feature of a mobile process calculus, and it is what makes the
calculus a specification of interaction rather than of a machine.

An implementation must close the gap. RSpace closes it by presenting the parallel
fragment as a store from channels to polarity-homogeneous bags and letting the
physical race between cores decide which produce is zipped with which consume.
That is a resolution mechanism, and it is a choice --- in the precise sense that
the family of admissible pairings, indexed by contended channels, has a section
only when something supplies one.

The observation this note develops is that the resolution mechanism is a
\emph{parameter}, and that other values of the parameter are available. In
particular, quantum mechanics is a resolution mechanism: it supplies amplitudes
for the branches of a race and the Born rule for turning them into outcomes.
Instantiating rho's parameter with that mechanism is a well-posed thing to do,
and the question this note answers is what syntax a programmer needs in order to
control it.

\subsection{The programme, in two stages}

The work proceeds in two stages, and it is worth separating them sharply because
they have different success criteria.

\emph{Stage one} exposes the additional structure. A quantum resolver
distinguishes programs that classical observation identifies, because it permits
interference. That is a failure of full abstraction --- the denotation is finer
than contextual equivalence --- and the failure is the feature. The set of pairs
that classical observation identifies and the quantum denotation separates is
the specification document for stage two.

\emph{Stage two} gives that structure syntactic expression, so that a programmer
can write into the gap on purpose. Aaronson's summary of quantum programming is
apt here: most of the discipline consists of arranging for interference to do
something useful. A syntax for arranging interference is therefore a syntax for
exploring what quantum advantage is available, and having one is a precondition
for a verdict rather than a consequence of one.

The weighted-GSLT programme proposed \emph{weight maps} for stage two: a table,
carried alongside the term, from spatial-behavioural formulae refining the
left-hand side of a rewrite rule to weights. This note argues that the weight map
had the right \emph{language} and the wrong \emph{location}, and that relocating
it into the for-comprehension's \texttt{where} clause --- by the simple device of
letting formulae take non-crisp truth values --- both subsumes it and repairs
three separate defects at once.

\subsection{Contributions}

\begin{enumerate}[leftmargin=1.4em,itemsep=3pt]
\item \textbf{A negative result about full abstraction} (\S\ref{sec:fa}).
Enlarging the class of rho contexts cannot expose interference, because contexts
can \emph{create} apertures but cannot \emph{read} them. The mismatch is in the
observation type, not the context class, and the correct frame is a
\emph{separation} between resolvers rather than a full-abstraction theorem
between terms. We give the replacement correctness condition: one-sided support
adequacy, whose failure to be two-sided is the interference.

\item \textbf{A sharp stage-one verdict} (\S\ref{sec:stage1}). Under the
race equation with uniform weights, unmodified rho interferes exactly when the
racing messages are structurally congruent, and then only constructively, because
the surviving message is a which-path record. The phenomenon is a permanent; it
is BosonSampling-shaped, hence plausibly hard and definitely not universal.

\item \textbf{The design} (\S\S\ref{sec:design}--\ref{sec:semantics}). Value the
\texttt{where} clause in a commutative semiring and evaluate it per candidate
match. Boolean values recover rholang 1.4 on the nose
(Theorem~\ref{thm:conservativity}).

\item \textbf{Weight maps are a normal form} (\S\ref{sec:weightmaps},
Theorem~\ref{thm:dnf}). Every weight map is a graded clause of the shape
$\bigoplus_i v_i \tensorv \crisp{\chi_i}$, and the partition discipline exists
solely to make that shape single-valued --- a requirement a formula satisfies for
free. Dynamic weight maps become formulae mentioning the current term, so the
update function disappears too (Proposition~\ref{prop:plastic}).

\item \textbf{The join is the recombiner} (\S\ref{sec:erasure},
Theorem~\ref{thm:scattering}). Which-path erasure --- the half of the problem the
weight map never addressed --- requires no new term former. A join
\texttt{\&} over $m$ patterns on a contended channel has $m!$ candidate matches
whose outcomes coincide precisely when the continuation is symmetric in the bound
names, and the graded clause supplies the $m!$ amplitudes. The clause is a
scattering matrix and the symmetric outcome carries its permanent.

\item \textbf{Shor's algorithm, and an overshoot} (\S\ref{sec:shor}). The
fragment is universal for quantum circuits (Lemmas~\ref{lem:1q}--\ref{lem:2q}),
the circuit's DAG is the term's causal order (Proposition~\ref{prop:dag}), and
the classical part costs nothing (Corollary~\ref{cor:modexp}). But with clauses
unconstrained the fragment decides $\mathsf{PP}$ by post-selection
(Proposition~\ref{prop:postbqp}), so the isometry condition of
Definition~\ref{def:isometric} is a containment requirement rather than a
convenience.

\item \textbf{PLN-valued clauses} (\S\ref{sec:pln}). The same slot accepts
Probabilistic Logic Network truth values, making inference and execution one
loop. We give the algebra honestly, including which of PLN's operations supply
values and which would need weakening to supply a semiring.

\item \textbf{A mortal forager} (\S\ref{sec:forager}). A second worked example,
in the inference-driven regime. The complementary guard pair's dependence on
negation in the hypothesis language dissolves
(Proposition~\ref{prop:trichotomy}); the power has a fixed point at zero, so a
belief needs prior evidence or an exploration disjunct
(Proposition~\ref{prop:bootstrap}); confidence is invisible without an outside
option (Remark~\ref{rem:cancel}); and the personality parameter $k$ acquires a
metabolic optimum (Principle~\ref{prin:k}) --- a return to PLN as well as from it.
\end{enumerate}

\subsection{What this note does not do}

It does not prove a quantum advantage; it builds the instrument with which one
could be sought. It does not settle the normalisation question raised in
\S\ref{ssec:fork}, though it makes the fork explicit and recommends a branch. It
does not treat replication or the full recursive fragment beyond noting where the
budget must bite. And it does not attempt a categorical account of the
construction, which we expect to be a Kleisli-style story over the semiring's
free-module monad and which we leave to a companion.

%%%%%%%%%%%%%%%%%%%%%%%%%%%%%%%%%%%%%%%%%%%%%%%%%%%%%%%%%%%%%%%%%%%%%%%%%%%%%%%
\section{Why context-based full abstraction cannot help}\label{sec:fa}

\subsection{An asymmetry between two translations}

When one encodes the $\pi$-calculus into rho, a natural worry about full
abstraction arises and a natural repair is available: rho contexts see strictly
more than $\pi$ contexts do, so demand preservation only against the images of
$\pi$ contexts. The repair is legitimate because the discrepancy it manages ---
rho's ability to inspect the structure of a name --- lies on an axis about which
\emph{both} calculi are silent in the same way, namely how a race is resolved.
Restricting the context class is orthogonal to resolution, so it costs nothing on
that axis.

The translation from rho into quantum mechanics is not like this. Quantum
mechanics makes a definite commitment about resolution: amplitudes compose
linearly and outcomes obey the Born rule. The discrepancy is therefore
\emph{on} the axis about which rho is silent. No restriction or enlargement of
the class of rho contexts can address it, for the simple reason that a rho
context is itself a rho term, whose own races are resolved by the same
unspecified mechanism. Composing silence with silence yields silence.

\begin{principle}[Contexts create apertures; they do not read them]\label{prin:aperture}
A rho context can change which races exist --- by adding contention, by
consuming a contender, by installing a new receive. It cannot observe how any
race is resolved. Consequently no property of the form ``for all rho contexts
$C[-]$, \dots'' can be sensitive to the resolution mechanism.
\end{principle}

This is the precise reason a full-abstraction theorem is not the right
correctness statement for the present setting, and the reason is stronger than
the usual observation that interference makes the denotation too fine. It is not
that we have too few contexts. It is that the discriminating power we want is not
of the kind contexts have.

\subsection{The door that stays open}

The enterprise is not thereby vacuous, and it is worth being careful about why.
A context cannot see a phase. But a context can be an \emph{interferometer}: an
arrangement in which a difference of phase becomes a difference in the
\emph{support} of the outcomes. This is how every physical quantum experiment
reports its result --- not by observing an amplitude but by a detector that never
fires. \S\ref{sec:erasure} exhibits such a context in three lines of rholang.

What this does not do is rescue full abstraction, because the discrimination it
achieves is between two \emph{resolvers} at a fixed term, not between two terms
at a fixed resolver.

\subsection{Separation as the correctness frame}

The right object of study is therefore the map
\[
  \textsf{Resolvers} \times \textsf{Terms} \longrightarrow \textsf{Distributions
  over normal forms},
\]
in which the classical RSpace semantics and the Born semantics are two points of
the first coordinate. Quantum advantage is the assertion that the image of the
quantum resolver is not contained in the image of the classical family --- a
separation, in the shape of a Bell argument, rather than a full-abstraction
theorem.

What survives as a checkable correctness obligation is one-sided.

\begin{definition}[Support adequacy and the interference deficit]\label{def:deficit}
Let $\Reach(P) \subseteq \NF$ be the set of normal forms operationally reachable
from $P$, and let $\sem{P}$ be the graded denotation of \S\ref{sec:semantics}.
The denotation is \emph{support-adequate} when
\[
  \supp(\sem{P}) \;\subseteq\; \Reach(P).
\]
The \emph{interference deficit} of $P$ is
$\;\mathcal{D}(P) := \Reach(P)\setminus\supp(\sem{P})$.
\end{definition}

Support adequacy says the interpretation invents nothing the operational
semantics forbids. The converse inclusion is exactly what destructive
interference violates: a normal form that some derivation reaches, whose
amplitudes cancel. So the natural two-sided condition is not merely unproven, it
is \emph{false precisely where the interesting physics is}, and its failure is
the observable. We will return to $\mathcal{D}$ as the quantity a programmer
manipulates.

\begin{remark}[Non-compositionality is not a defect]\label{rem:compositional}
The store denotation of rho is a monoid homomorphism for parallel composition,
because it denotes the store. The sum-over-futures denotation cannot be, because
composing $P$ with $Q$ creates races that neither had alone. This is not a
weakness relative to quantum mechanics; it is the same phenomenon in the same
place --- interacting subsystems do not have product states. The cut is the
interaction term, and entanglement across a cut is the semantic image of a race
across a cut. It is also the semantic shadow of the observation that an aperture
is a property of the cut and not of either agent flanking it.
\end{remark}

%%%%%%%%%%%%%%%%%%%%%%%%%%%%%%%%%%%%%%%%%%%%%%%%%%%%%%%%%%%%%%%%%%%%%%%%%%%%%%%
\section{What unmodified rho already gives}\label{sec:stage1}

\subsection{The race equation}

Take as the constraint on an interpretation $\sem{-}$, at a two-way race between
one receive and two sends,
\begin{equation}\label{eq:race}
  \sem{\,x!(Q_1) \mid \textsf{for}(y \Leftarrow x)P \mid x!(Q_2)\,}
  \;=\;
  w_1\,\sem{\,P\{@Q_1/y\} \mid x!(Q_2)\,}
  \;+\;
  w_2\,\sem{\,x!(Q_1) \mid P\{@Q_2/y\}\,},
\end{equation}
with $w_1 = w_2$ and the weights covering the total probability space; dually for
two receives racing over one message. The constraint says nothing about a lone
receive meeting a lone send. Alongside it we adopt the principle that the meaning
of a term is the sum of all its possible futures.

Two things are worth noticing immediately about \eqref{eq:race}. It is
\emph{parameter-free}: the weights are forced by the arity of the race, so there
is no table, no key language, and no partition. And it is \emph{silent on the
confluent fragment}, which localises whatever design freedom remains.

\subsection{The residual is a which-path record}

The left summand of \eqref{eq:race} retains $x!(Q_2)$ and the right retains
$x!(Q_1)$. That surviving message records which branch was taken. Two derivations
whose residuals differ do not reconverge, and derivations that do not reconverge
cannot interfere.

There is exactly one escape, and it is forced: when $Q_1 \scong Q_2$ the record
is vacuous, the two summands are structurally congruent, and the branches meet.

\begin{proposition}[Stage-one verdict]\label{prop:stage1}
Under \eqref{eq:race} with uniform real weights and no further apparatus, the
only interference available to rho is between derivations consuming
structurally congruent messages, and it is always constructive.
\end{proposition}

\begin{proof}[Proof sketch]
Interference requires two derivations from a common source reaching a common
target. At a single race the targets are $P\{@Q_1/y\} \mid x!(Q_2)$ and
$x!(Q_1) \mid P\{@Q_2/y\}$, which are congruent only if $Q_1 \scong Q_2$ (the
substitution instances agreeing for all $P$ forces the quoted payloads to agree).
Iterating, any reconvergence decomposes into steps of this form together with
permutations of independent redexes, which are identified by
Definition~\ref{def:causal} below. All weights being non-negative reals, summed
amplitudes cannot cancel.
\end{proof}

\begin{remark}[The general race is a permanent]\label{rem:permanent}
At a channel carrying a bag of produces and a bag of consumes, communication
chooses an admissible pairing and zips, leaving surplus. Summing coherently over
perfect matchings of an amplitude matrix is the computation of a permanent, which
is to say that the coherent reading of the zip is BosonSampling-shaped: plausibly
hard to sample, while the classical reading --- pick a matching uniformly ---
is easy. The superlinear multiplicity that earlier work identified in the
degenerate case is accordingly not a bookkeeping error but genuine
indistinguishable-particle statistics.
\end{remark}

The verdict is therefore genuinely good news for the design, because it makes the
brief precise: what rho lacks is not interference \emph{per se} but \emph{phase}
and \emph{control over recombination}. The analogy with linear optics is close
enough to be useful. Indistinguishable bosons through a passive network give
hardness without universality; adding phase control and adaptivity is what
carries one to KLM-style universality. The gap to be closed may be small.

\subsection{The normalisation fork}\label{ssec:fork}

Equation~\eqref{eq:race} is not consistent as stated, and the inconsistency shows
up at its own simplest instance. With $Q_1 \scong Q_2$ the right-hand side is
$(w_1 + w_2)\sem{T}$ for a single $T$. If the $w_i$ are amplitudes with
$|w_1|^2 + |w_2|^2 = 1$, this has norm $\sqrt{2}$; for an $m$-fold degenerate
race, norm $\sqrt{m}$ and weight $m$.

Three resolutions are available.

\begin{enumerate}[leftmargin=1.4em,itemsep=2pt]
\item \emph{Convex.} Let $\sem{-}$ take values in formal convex combinations,
$w_i = 1/m$. Consistent, and exactly the classical scheduler with a fair coin.
No interference.
\item \emph{Branch register.} Adjoin the derivation index so the summands stay
distinguishable. Consistent and interference-capable, at the cost of carrying the
index --- this is the route taken by the ZX compilation work.
\item \emph{Unnormalised.} Let $\sem{-}$ take values in the free $\V$-module on
normal forms and apply the Born rule at observation. Consistent, with $m$-fold
degeneracy carrying weight $m^2$.
\end{enumerate}

We take the third, on the ground that linear-in-$m$ and quadratic-in-$m$ are the
values of two different observables rather than two normalisations of one: the
first is what one gets by looking at which derivation fired, the second by
declining to look. Since the whole point of the exercise is to decline to look,
the quadratic answer is the correct one here, and it is exactly the bosonic
enhancement of Remark~\ref{rem:permanent}.

%%%%%%%%%%%%%%%%%%%%%%%%%%%%%%%%%%%%%%%%%%%%%%%%%%%%%%%%%%%%%%%%%%%%%%%%%%%%%%%
\section{Two requirements, derived}\label{sec:requirements}

Interference needs two things and nothing else: derivations that
\emph{reconverge}, and a nonzero \emph{relative phase} between them. Every
syntactic feature proposed below is answerable to one of these.

\subsection{Reconvergence must be counted over causal histories}

A caution first, because it disciplines the path sum. Every concurrency diamond
is a reconvergence: two independent redexes fired in either order reach the same
term. Rho programs are saturated with these. If the sum over futures ranges over
\emph{interleavings}, then $m$ independent redexes contribute a factor of $m!$
arising from nothing but the number of linearisations of a partial order.

\begin{definition}[Independence and causal histories]\label{def:causal}
Two redex firings at a term are \emph{independent} when their footprints ---
the data and continuations they consume --- are disjoint. Let $\sim$ be the least
congruence on derivations identifying $u\,;\,v$ with $v\,;\,u$ for independent
$u,v$. A \emph{causal history} is a $\sim$-class of derivations, and the path sum
of \S\ref{ssec:pathsum} ranges over causal histories.
\end{definition}

\begin{proposition}[No spurious combinatorial enhancement]\label{prop:causal}
Summing over causal histories rather than interleavings removes the $m!$ factor,
and does not remove any reconvergence between causally inequivalent derivations.
\end{proposition}

The apparatus for this choice already exists in the history-monad account of
rho: free histories form the universal cover of the reduction graph and the
erasure lattice is indexed by subgroups of its fundamental group. Filling the
concurrency squares is one such erasure, and interference lives on the loops
that survive it.

\subsection{The two syntactic requirements}

\begin{enumerate}[label=(R\arabic*),leftmargin=2.6em,itemsep=3pt]
\item \label{r1} \textbf{Phase conditional on data.} The amplitude attaching to a
branch must be allowed to depend on which datum that branch consumed, and on the
structure of the term around it.
\item \label{r2} \textbf{Which-path erasure.} The programmer must be able to
arrange that two branches produce the \emph{same} residual, so that their
amplitudes are added rather than kept apart.
\end{enumerate}

\begin{remark}[Auditing weight maps against the requirements]\label{rem:audit}
A weight map answers \ref{r1} well: its keys are spatial-behavioural formulae
refining the left-hand side of a rule, which is precisely ``a phase conditional
on the structure of what was matched.'' It says nothing whatever about \ref{r2};
in the ZX treatment, recombination arrives as a choice of measurement basis at
the end, which is an observer's decision made outside the program. And it is
ergonomically misplaced: a table carried in the state is a Hamiltonian
specification, whereas arranging interference is done with term formers, at the
site where the branching happens.
\end{remark}

%%%%%%%%%%%%%%%%%%%%%%%%%%%%%%%%%%%%%%%%%%%%%%%%%%%%%%%%%%%%%%%%%%%%%%%%%%%%%%%
\section{The design}\label{sec:design}

The design is one change, stated in a sentence: \emph{the \texttt{where} clause
of the for-comprehension is valued in a commutative semiring and is evaluated
once per candidate match.}

\subsection{Resolution algebras}

\begin{definition}[Resolution algebra]\label{def:algebra}
A \emph{resolution algebra} is a commutative semiring
$\V = (V, \sumv, \tensorv, 0, 1)$: $\sumv$ associative and commutative with unit
$0$, $\tensorv$ associative and commutative with unit $1$, $\tensorv$
distributing over $\sumv$, and $0$ annihilating.
\end{definition}

The two operations are forced by what the connectives must mean.

\begin{itemize}[leftmargin=1.4em,itemsep=2pt]
\item \emph{Conjunction is $\tensorv$.} Two conditions holding together compose
multiplicatively; this is also what makes the joined receive
$p_1 \Leftarrow c_1\ \&\ p_2 \Leftarrow c_2$ combine its clauses correctly.
\item \emph{Disjunction is $\sumv$.} Two ways for a match to be enabled add.
\end{itemize}

\begin{principle}[Interference is a logical phenomenon]\label{prin:or}
Since $\sem{\varphi \vee \psi} = \sem{\varphi} \sumv \sem{\psi}$, in a resolution
algebra with additive inverses a disjunction may evaluate to $0$ though neither
disjunct does. Destructive interference is a \emph{false disjunction with true
disjuncts}. This is the whole of the additional structure, stated in the logic
rather than imposed on it from outside.
\end{principle}

\begin{table}[t]
\centering
\small
\begin{tabular}{@{}llll@{}}
\toprule
\textbf{Algebra} & \textbf{$\sumv,\tensorv$} & \textbf{Resolver} & \textbf{Regime}\\
\midrule
$\Bool = (\{0,1\},\vee,\wedge)$ & $\vee$, $\wedge$ & demonic / RSpace race & rholang 1.4 exactly \\
$\Rnn = (\mathbb{R}_{\geq 0},+,\times)$ & $+$, $\times$ & Gillespie / SSA & stochastic \\
$\Cx = (\mathbb{C},+,\times)$ & $+$, $\times$ & Born rule & interference \\
Viterbi $([0,1],\max,\times)$ & $\max$, $\times$ & greedy & most-likely path \\
Tropical $(\mathbb{R}\cup\{\infty\},\min,+)$ & $\min$, $+$ & least cost & cost-optimal path \\
$\pow$-PLN $([0,1],+,\times)$ & $+$, $\times$ & inference-driven & \S\ref{sec:pln} \\
\bottomrule
\end{tabular}
\caption{Instances. One syntax; the resolver is the choice of $\V$.}\label{tab:algebras}
\end{table}

Negation is the casualty and it is worth saying so plainly. On $[0,1]$ the map
$v \mapsto 1 - v$ serves; on $\Cx$ there is no canonical complement and De
Morgan fails against $(+,\times)$. The design therefore keeps two sorts, with the
crisp sort embedded in the graded one.

\subsection{Two sorts of formula}

Let the crisp sort be the existing condition language of rholang 1.4: ground
formulae supplied by a language definition, the spatial connectives generated
from the term formers, the context-labelled modality $\langle K \rangle$, and
full Boolean structure.
\begin{align*}
\beta &::= \textsf{true} \mid \textsf{false} \mid g \mid \neg\beta \mid
  \beta \wedge \beta \mid \beta \vee \beta \mid \langle K \rangle\beta \mid
  \textsf{out}(n)\beta \mid \textsf{in}(n)\beta \mid \beta \mid \beta \mid @\beta
\intertext{and let the graded sort be}
\psi &::= \crisp{\beta} \mid v \mid \psi \tensorv \psi \mid \psi \sumv \psi
  \mid \langle K \rangle_{\V}\,\psi \qquad (v \in V).
\end{align*}
Here $\crisp{\beta}$ is the indicator embedding, sending a satisfied $\beta$ to
$1$ and an unsatisfied one to $0$, and $v$ ranges over scalar constants of the
algebra. The graded sort has no negation and no fixed-point operator; the crisp
sort keeps both.

\begin{remark}[One slot, not two]\label{rem:oneslot}
It is tempting to give the for-comprehension two clauses, a Boolean gate and a
graded weight. It is better to give it one graded clause and recover the gate as
its \emph{support}: a communication is enabled exactly when its clause evaluates
to something other than $0$. This keeps the language unforked --- existing
Boolean where-clauses are the $\Bool$-valued fragment under
$\Bool \hookrightarrow \V$ --- and it makes support adequacy
(Definition~\ref{def:deficit}) hold by construction rather than by proof, since
the classical interpreter sees nothing of the value except whether it vanishes.
\end{remark}

\subsection{Syntax}

The for-comprehension of rholang 1.4 already carries a \texttt{where} clause,
whose condition may mention the variables bound by the preceding receipts. The
change is only to the sort of that condition.

\begin{lstlisting}
for( p11 <- c11 & ... & pm1 <- cm1 ;
     ... ;
     p1n <- c1n & ... & pmn <- cmn
     where psi ) P
\end{lstlisting}

\noindent
with \texttt{psi} a graded formula over the algebra $\V$ declared for the shard
or the channel's language definition. Within \texttt{psi}, the bound names
$p_{ij}$ are in scope, so the value may depend on what was matched --- which is
requirement \ref{r1}.

\subsection{Per-match evaluation}\label{ssec:permatch}

The decisive point, and the one from which everything in \S\ref{sec:erasure}
follows, concerns \emph{when} the clause is evaluated.

\begin{definition}[Candidates]\label{def:cand}
Let $r$ be a receipt site in $P$ --- a for-comprehension together with the
channels its receipts name. A \emph{candidate} at $r$ is an assignment $\sigma$
of a distinct datum occurrence to each pattern of each group, such that every
pattern matches its assigned datum. Write $\Cand(P,r)$ for the set of candidates
and $\Out(r,\sigma) \in \Proc/{\scong}$ for the residual term the firing of $r$
under $\sigma$ produces.
\end{definition}

Candidates are indexed by \emph{occurrences}, not by terms. This matters: two
structurally congruent messages on a channel are two candidates, not one,
because parallel composition is a multiset and structural congruence does not
delete duplicates. An implementation that keys data by content alone, or that
allocates identifiers during firing, cannot represent $\Cand$ faithfully.

\begin{definition}[Graded clause valuation]\label{def:val}
A graded clause $\psi$ at a receipt site $r$ denotes a function
\[
  \sem{\psi}_r : \Cand(P,r) \longrightarrow V ,
\]
defined by structural recursion, with $\crisp{\beta}$ evaluated against the term
$P$ under the substitution $\sigma$ carries, scalars constant, $\tensorv$ and
$\sumv$ pointwise, and $\langle K \rangle_\V$ as in \S\ref{ssec:modality}.
\end{definition}

So a clause is a \emph{vector} indexed by candidates, and --- once we compose it
with $\Out$ --- a \emph{matrix} over candidates and outcomes. That is the whole
mechanism.

\subsection{The graded modality is a transfer operator}\label{ssec:modality}

Let $\V\langle \Proc/{\scong}\rangle$ be the free $\V$-module on terms modulo
structural congruence. Define the one-step operator of a term by summing its
candidates:
\begin{equation}\label{eq:onestep}
  T(P) \;=\; \sum_{r}\ \sum_{\sigma \in \Cand(P,r)}
    \sem{\psi_r}_r(\sigma)\;\cdot\;\Out(r,\sigma) ,
\end{equation}
the outer sum over receipt sites, taken over causal histories in the sense of
Definition~\ref{def:causal}, and the coefficient arithmetic in $\V$. Then the
graded modality is
\[
  \sem{\langle K \rangle_\V\,\psi}(P)
  \;=\; \sum_{r,\sigma} \sem{\psi_r}_r(\sigma) \tensorv \sem{\psi}(\Out(r,\sigma)),
\]
which is $T$ acting on the vector of values of $\psi$. In other words the graded
modality is the transfer matrix of the reduction relation, and the principle that
the meaning of a term is the sum of all its possible futures is exactly the
statement that $\sem{-}$ is a fixed point of $\langle K\rangle_\V$. The
denotational principle and the graded logic are one object seen twice.

%%%%%%%%%%%%%%%%%%%%%%%%%%%%%%%%%%%%%%%%%%%%%%%%%%%%%%%%%%%%%%%%%%%%%%%%%%%%%%%
\section{Formal semantics}\label{sec:semantics}

\subsection{Terms}

We work with core rho, taken with a built-in name restriction whose desugaring is
not our concern here.
\[
  P ::= 0 \;\mid\; x!(P) \;\mid\; \textsf{for}(\,\vec{y} \Leftarrow \vec{x}
  \ \textsf{where}\ \psi\,)P \;\mid\; P \mid P \;\mid\; *x \;\mid\; \textsf{new}\ x\ \textsf{in}\ P,
  \qquad
  x ::= @P .
\]
Quoting is $@P$ and dropping is $*x$. Structural congruence $\scong$ is the least
congruence making $(\Proc, \mid, 0)$ a commutative monoid, together with the
usual laws for \textsf{new} and $*@P \scong P$.

\subsection{The state space}

\begin{definition}[Graded denotation]\label{def:denotation}
Fix a resolution algebra $\V$ and a budget $\beta \in \Nat$. Let
$M := \V\langle \Proc/{\scong}\rangle$ be the free $\V$-module on terms modulo
structural congruence, with basis vectors written $\lvert P \rangle$. The
$\beta$-budgeted denotation is
\[
  \sem{P}_\beta \;:=\; T_{\!*}^{\,\beta}\,\lvert P \rangle
  \;=\;\!\!\sum_{h \,:\, P \Rightarrow^{\leq\beta} Q}\!\! w(h)\,\lvert Q\rangle ,
\]
where $T_*$ extends $T$ of \eqref{eq:onestep} linearly, $h$ ranges over causal
histories of length at most $\beta$ ending at a normal form or exhausting the
budget, and $w(h)$ is the $\tensorv$-product of the clause values along $h$.
\end{definition}

\begin{remark}[Why a budget, and only one]\label{rem:budget}
The definition is corecursive and does not bottom out on non-terminating terms.
For terminating $P$ the recursion is well-founded on the reduction tree and
$\beta$ may be taken to be its depth. In general one needs either a convergence
hypothesis --- a spectral radius condition on $T$, since $\Cx$ has no order and
Knaster--Tarski is unavailable --- or a budget. We take the budget, and note it
is the same budget already charged elsewhere: the risk ledger's single $\omega$
for reflection, the address-space bound of the compilation route, and the
inference-token budget that keeps generated model-checking terminating. These
should be one budget doing three jobs rather than three budgets.
\end{remark}

To make normal forms absorbing and keep histories comparable in length we adopt
the usual convention.

\begin{definition}[Fictitious jump]\label{def:fict}
At a normal form $Q$ the operator acts as $T\lvert Q\rangle = \lvert Q\rangle$.
\end{definition}

\subsection{Readings}\label{ssec:pathsum}

One denotation admits three readings, differing only in what is done at the end.

\begin{enumerate}[leftmargin=1.4em,itemsep=2pt]
\item \emph{Support.} $\supp(\sem{P}_\beta) \subseteq \NF$. This is what the
classical interpreter sees, and by Remark~\ref{rem:oneslot} it is all it sees.
\item \emph{Born.} For $\V = \Cx$, the distribution
$Q \mapsto |\langle Q \mid \sem{P}_\beta\rangle|^2$, normalised.
\item \emph{Stochastic.} For $\V = \Rnn$, the same expression without the
square, which is the jump-chain marginal of the Gillespie reading.
\end{enumerate}

\begin{theorem}[Conservativity]\label{thm:conservativity}
Let $\V = \Bool$ and let every clause be the embedding $\crisp{\beta}$ of a
crisp condition. Then $\supp(\sem{P}_\beta)$ is exactly the set of normal forms
reachable from $P$ in at most $\beta$ steps under rholang 1.4's operational
semantics with its Boolean \texttt{where} clauses.
\end{theorem}

\begin{proof}[Proof sketch]
In $\Bool$, $\tensorv = \wedge$ and $\sumv = \vee$, so $w(h) = 1$ iff every step
of $h$ has a satisfied clause, and the coefficient of $\lvert Q\rangle$ is the
disjunction over histories reaching $Q$. Candidates coincide with the admissible
zips of the store presentation, and the fictitious jump does not add normal
forms. Hence the coefficient is $1$ iff some admissible derivation reaches $Q$.
\end{proof}

\begin{corollary}[Support adequacy]\label{cor:adequacy}
For any $\V$, $\supp(\sem{P}_\beta) \subseteq \Reach_\beta(P)$, since a nonzero
coefficient requires at least one history with all clause values nonzero, and
such a history is admissible under the Boolean image of the clauses. Equality
fails exactly when some $\Reach$-witnessed coefficient is a sum of nonzero terms
summing to $0$ --- that is, exactly on the interference deficit
$\mathcal{D}(P)$.
\end{corollary}

This is the payoff of Remark~\ref{rem:oneslot}. Adding grading cannot make a
classical program do something it could not do; it can only make it fail to do
something it could.

%%%%%%%%%%%%%%%%%%%%%%%%%%%%%%%%%%%%%%%%%%%%%%%%%%%%%%%%%%%%%%%%%%%%%%%%%%%%%%%
\section{Weight maps are a normal form}\label{sec:weightmaps}

\begin{definition}[Weight map, recalled]\label{def:wmap}
A weight map for a rewrite rule $R$ is a finite family of spatial-behavioural
formulae $\chi_1,\dots,\chi_k$ refining the left-hand side of $R$, required to
\emph{partition} that refinement, together with values $v_1,\dots,v_k \in V$;
the weight of a redex is $v_i$ for the unique $i$ with $\chi_i$ satisfied.
\end{definition}

\begin{theorem}[Weight maps are the linear graded clauses]\label{thm:dnf}
Every weight map $(\chi_i,v_i)_{i\leq k}$ is denoted by the graded clause
\[
  \psi \;=\; \bigoplus_{i=1}^{k} \; v_i \tensorv \crisp{\chi_i},
\]
and conversely a graded clause of this shape denotes a weight map if and only if
the $\chi_i$ partition the refinement. Graded clauses not of this shape ---
those with nested $\tensorv$ over overlapping conditions, or with a graded
modality --- are strictly more expressive.
\end{theorem}

\begin{proof}[Proof sketch]
Forwards: on a candidate satisfying $\chi_j$ and no other, all but the $j$th
summand vanish and the value is $v_j$, and totality of the partition ensures some
summand survives. Backwards: if two $\chi$'s overlap on a candidate, the clause
returns $v_i \sumv v_j$, which is a well-defined value but not a table lookup;
the partition condition is exactly the condition making the sum a selection.
Strictness: with overlapping $\chi_i$ the clause is a genuine mixture no weight
map assigns, and $\langle K\rangle_\V$ inspects successors, which no key can.
\end{proof}

\begin{corollary}[The partition discipline is unnecessary]\label{cor:nopartition}
Single-valuedness of a graded clause is automatic, a formula having one value.
The requirements that key sets partition the left-hand-side refinement, that
partitions be constructed through a single checked funnel, and that a default
class complete them, all exist to secure single-valuedness for a lookup and have
no analogue here.
\end{corollary}

This also dissolves a specific obstruction. Earlier work observed a tension
between the \emph{Boolean demand} --- a weight map's partition needs complements,
so the target logic specification must supply negation --- and the fact that the
transport of a generated logic along a bisimulation-preserving translation
preserves only the geometric fragment, which has no complements. With no
partition to define, the Boolean demand does not arise. Negation is still absent
from the graded sort, but it is absent as a design choice about the value algebra
rather than as an unmet requirement.

\begin{proposition}[Plasticity is free]\label{prop:plastic}
A dynamic weight map, whose table is revised after each rewrite by an update
function $u : W \times \textsf{Match} \times \textsf{Trace} \to W$, is denoted by
a static graded clause whenever $u$'s dependence on the trace is expressible as a
crisp formula over the current term. No update function is needed, because a
clause is re-evaluated against the current term at every candidate.
\end{proposition}

\begin{proof}[Proof sketch]
The clause $\psi$ mentions the term through its crisp subformulae; its value at
a candidate in the term $P'$ reached after some history is by definition computed
against $P'$. A table revised to depend on features of $P'$ therefore agrees with
the clause $\bigoplus_i v_i \tensorv \crisp{\chi_i}$ in which the $\chi_i$
mention those features directly. Dependence on trace history beyond what is
visible in $P'$ is not expressible, which is the hypothesis.
\end{proof}

\begin{remark}
Proposition~\ref{prop:plastic} retires a specific implementation obligation. The
plan for compiling the weighted simulator identified the update function --- an
arbitrary closure over time and a trace summary --- as the one component with no
compilable form, and proposed a declarative sub-language for it. The graded
clause \emph{is} that sub-language, and it was already needed for other reasons.
The restriction to features visible in the current term has precedent: the
exhaustive-graph construction already refuses time-dependent updates by fixing
the clock at zero.
\end{remark}

%%%%%%%%%%%%%%%%%%%%%%%%%%%%%%%%%%%%%%%%%%%%%%%%%%%%%%%%%%%%%%%%%%%%%%%%%%%%%%%
\section{The join is the recombiner}\label{sec:erasure}

Requirement \ref{r2} asked for which-path erasure, and the audit of
Remark~\ref{rem:audit} found nothing in the weight map that supplies it. It turns
out no new term former is needed: the erasure is already in the language, in the
join.

\subsection{Permutations of a join are the interfering pair}

Consider a receipt whose group contains $m$ patterns on a single contended
channel, meeting $m$ data occurrences:
\[
  x!(Q_1) \mid \cdots \mid x!(Q_m) \mid
  \textsf{for}(\,y_1 \Leftarrow x\ \&\ \cdots\ \&\ y_m \Leftarrow x
  \ \textsf{where}\ \psi\,)P .
\]
The candidates are the $m!$ bijections $\sigma$ from patterns to data. Every
candidate consumes \emph{all} the data --- there is no surplus, hence no residual
message and no which-path record --- and produces
$\Out(\sigma) = P\{@Q_{\sigma(1)}/y_1,\dots,@Q_{\sigma(m)}/y_m\}$.

\begin{definition}[Symmetric continuation]\label{def:sym}
$P$ is \emph{symmetric} in $y_1,\dots,y_m$ when
$P\{@Q_{\sigma(1)}/y_1,\dots\} \scong P\{@Q_{\tau(1)}/y_1,\dots\}$ for all
bijections $\sigma,\tau$ and all payloads.
\end{definition}

Symmetry is cheap to arrange, because parallel composition is commutative: the
continuation $z!(*y_1) \mid z!(*y_2)$ is symmetric, while
$z_1!(*y_1) \mid z_2!(*y_2)$ is not.

\begin{theorem}[The clause is a scattering matrix]\label{thm:scattering}
Let $A_{ij} := \sem{\psi}(\,y_i \mapsto @Q_j\,)$ be the value of the clause on the
candidate assigning datum $j$ to pattern $i$, and suppose $\psi$ factors as
$\bigotimes_i \psi_i(y_i)$ over patterns. If $P$ is symmetric in the $y_i$, all
$m!$ candidates share a single outcome $\bar{P}$, and
\[
  \langle \bar{P} \mid T \lvert P\text{-site}\rangle
  \;=\; \sum_{\sigma \in S_m} \prod_{i=1}^{m} A_{i\sigma(i)}
  \;=\; \Per(A).
\]
If $P$ is not symmetric, the candidates have pairwise distinct outcomes and each
carries its own product with no summation.
\end{theorem}

\begin{proof}[Proof sketch]
Every candidate is a bijection and the clause factors, so the weight of the
candidate $\sigma$ is $\prod_i A_{i\sigma(i)}$. Under symmetry all outcomes are
one basis vector and the coefficients add, which is the permanent. Without
symmetry the outcomes are distinct basis vectors and no addition occurs.
\end{proof}

\begin{corollary}[Which-path, syntactically]\label{cor:whichpath}
No choice of clause produces interference at a receipt whose continuation
distinguishes its bound names. Whether a race interferes is decided by the
continuation, not by the weights.
\end{corollary}

This is the design's central ergonomic fact. The programmer has two dials at
exactly the site where the branching happens: the clause sets the amplitudes, and
the continuation decides whether the branches are told apart. Arranging
interference is arranging for the continuation to forget.

\begin{remark}[Statistics]\label{rem:stats}
$\Per(A)$ is the bosonic amplitude. A clause assigning the sign of $\sigma$ ---
expressible when the algebra has additive inverses and the crisp sort can
detect the relevant order --- gives $\Det(A)$ and hence fermionic statistics, in
which coincident data annihilate. That the same syntactic slot produces both is a
point in favour of the slot.
\end{remark}

\subsection{The minimal interferometer}

Take $\V = \Cx$, $m = 2$, and a clause that returns $a$ on the identity
assignment and $b$ on the transposition. With a symmetric continuation:

\begin{lstlisting}
new x, z in {
  x!(Q1) | x!(Q2) |
  for( y1 <- x & y2 <- x where psi ) { z!(*y1) | z!(*y2) }
}
\end{lstlisting}

The single outcome $z!(*Q_1) \mid z!(*Q_2)$ carries amplitude $a + b$. Choosing
$b = -a$ makes it zero. That normal form is operationally reachable --- indeed it
is the \emph{only} reachable normal form --- so
$\mathcal{D}(P) \neq \emptyset$ and Corollary~\ref{cor:adequacy}'s inclusion is
strict. Replacing the continuation by $z_1!(*y_1) \mid z_2!(*y_2)$ restores two
distinct outcomes with amplitudes $a$ and $b$, probabilities $|a|^2$ and $|b|^2$,
and no cancellation whatever the phases.

Three lines of rholang thus exhibit the entire phenomenon: a detector that never
fires, and a which-path measurement that destroys the effect. It is the
Hong--Ou--Mandel arrangement, with the channel $z$ playing the role of a single
detector and the pair $z_1,z_2$ of two.

\begin{remark}[What a rho context can and cannot do]\label{rem:back-to-fa}
Note the consistency with \S\ref{sec:fa}. The context
$\textsf{for}(y_1 \Leftarrow x\ \&\ y_2 \Leftarrow x \dots)\{z!(*y_1)\mid z!(*y_2)\}$
does not observe the resolution of the race; it \emph{arranges} the race so that
resolution leaves a trace in the support. This is precisely what was meant by a
context acting as an interferometer, and it is the reason the enterprise has
observable consequences without any context ever reading an amplitude.
\end{remark}

%%%%%%%%%%%%%%%%%%%%%%%%%%%%%%%%%%%%%%%%%%%%%%%%%%%%%%%%%%%%%%%%%%%%%%%%%%%%%%%
\section{A worked example: Shor's algorithm}\label{sec:shor}

The design of \S\ref{sec:design} was derived from what interference requires,
not from the circuit model, so it is a fair question whether it reaches the
standard algorithms. It does. The instructive part is not that the fragment turns
out expressive enough but that it turns out \emph{too} expressive, which
\S\ref{ssec:hazard} addresses and which changes the status of one of this note's
open threads.

Every number reported in this section is computed by a simulator implementing
Definition~\ref{def:denotation} literally: it enumerates the candidates at each
site and adds coefficients on coinciding outcomes. Nothing in it multiplies gate
matrices. Matrices appear only as the values a clause returns on a candidate.

\subsection{A qubit and a gate}\label{ssec:gadget}

\begin{definition}[Wire encoding]\label{def:wire}
A \emph{qubit} is a channel carrying exactly one of two designated data. Writing
$a_0,a_1$ for the two, the basis states are $q!(a_0)$ and $q!(a_1)$.
\end{definition}

\begin{lstlisting}
new q, w, q' in {
  q!(*in) | w!(a0) | w!(a1) |
  for( v <- q & u1 <- w & u2 <- w where psi ) { q'!(*u1) }
}
\end{lstlisting}

\begin{lemma}[Arbitrary one-qubit gate]\label{lem:1q}
The site above has exactly two candidates, the bijections
$\{u_1,u_2\} \to \{a_0,a_1\}$. Each consumes every datum on $q$ and on $w$, so
no residual records which fired. The outcomes are $q'!(a_0)$ and $q'!(a_1)$;
they are distinct, and neither mentions $v$, hence they are the \emph{same two
terms} for either input encoding. Therefore the one-step operator restricted to
$\mathrm{span}\{q!(a_0), q!(a_1)\}$ is the matrix
$A_{ij} = \sem{\psi}(v \mapsto a_i,\; u_1 \mapsto a_j)$, and since $\psi$ may
mention both $v$ and $u_1$, every $2\times 2$ matrix over $\V$ is realised.
\end{lemma}

\begin{remark}[Two erasures, not one]\label{rem:twoerasures}
Theorem~\ref{thm:scattering} identified outcomes of \emph{different candidates at
one site}, which is recombination within a superposition. A gate needs the other
identification: outcomes at \emph{different input basis states} must coincide,
which is what makes the site a linear map into a fixed codomain rather than a
family of unrelated branchings. Both are bought by the same device --- a join
consumes everything, and the continuation re-emits without mentioning what was
consumed. That the join's totality purchases both erasures is the reason no
further term former is needed here either.
\end{remark}

\begin{lemma}[Arbitrary two-qubit gate]\label{lem:2q}
Joining two input channels with two ancilla channels gives
$2 \times 2 = 4$ candidates and four distinct outcomes indexed by the pair of
$u_1$-assignments. The clause may mention both input data and both assignments,
so every $4 \times 4$ matrix is realised, including entangling ones. Diagonal
gates are the case where the clause vanishes off the diagonal, and by
Remark~\ref{rem:oneslot} those candidates simply do not fire.
\end{lemma}

\subsection{Composition, and where the circuit's wires went}

\begin{proposition}[The circuit's DAG is the term's causal order]\label{prop:dag}
A gadget is enabled exactly when its input channels carry data. Composing
gadgets in series on fresh channels therefore makes the data dependency of the
term equal to the wire structure of the circuit, and every reduction order is a
linear extension of the circuit's dependency DAG. No scheduling discipline is
required and none is expressible: the term does not encode the circuit's order,
it \emph{is} that order.
\end{proposition}

Proposition~\ref{prop:causal} is load-bearing here rather than hygienic. Gates in
one circuit layer are independent redexes, so without the causal quotient a layer
of $n$ gates would contribute a factor $n!$ from counting linearisations.
Measured, for $n$ independent one-qubit gadgets: the causal sum gives squared
norm $1.000000$ for $n = 2,3,4$, while summing over interleavings inflates it by
exactly $2$, $6$ and $24$.

\subsection{The classical part is free}

\begin{proposition}[Deterministic communication is amplitude-transparent]\label{prop:free}
A communication with a single candidate contributes the clause's value on that
candidate, which for a crisp clause is $1$. Hence amplitudes are unchanged by
deterministic reduction, and the \emph{length} of a derivation does not affect
its weight.
\end{proposition}

\begin{corollary}\label{cor:modexp}
Any ordinary rho program runs inside every branch of a superposition at no cost
in the semantics. In particular modular exponentiation --- the expensive part of
Shor's algorithm, and the part usually requiring laborious reversible
compilation --- is written as ordinary rholang.
\end{corollary}

Measured: the block computing $x \mapsto 7^x \bmod 15$ carries a $16$-term
superposition to a $16$-term superposition with norm $1.000000000000$, exactly
preserved.

\subsection{The instance, assembled}

\begin{theorem}[Expressibility]\label{thm:shor}
For $N$ with $n = \lceil \log_2 N\rceil$, Shor's algorithm is a graded rho term
using $O(n^2)$ gadgets of arity at most two for the quantum Fourier transform,
$O(n^3)$ deterministic communications for modular exponentiation, and a term
size and reduction depth polynomial in $n$.
\end{theorem}

Two verifications.

\emph{The transform.} The QFT assembled from the gadgets of
Lemmas~\ref{lem:1q}--\ref{lem:2q} --- Hadamards, controlled phases and swaps, in
the standard decomposition --- was compared entry by entry against the discrete
Fourier matrix $\omega^{jk}/\sqrt{2^n}$. Worst entry error: $2.7\times 10^{-16}$
at $n=2$, $1.4\times 10^{-15}$ at $n=3$, $3.8\times 10^{-15}$ at $n=4$.

\emph{The instance.} $N=15$, $a=7$ (so $r=4$), with four counting and four work
qubits. The run fires $8$ one-qubit and $8$ two-qubit gadgets, evaluating $390$
candidates in all, and ends with $16$ configurations carrying nonzero amplitude.
The counting register reads

\begin{table}[h]
\centering
\small
\begin{tabular}{@{}lccccl@{}}
\toprule
outcome $c$ & $0$ & $4$ & $8$ & $12$ & all others \\
\midrule
probability & $0.250000000$ & $0.250000000$ & $0.250000000$ & $0.250000000$ & $0$ \\
\bottomrule
\end{tabular}
\caption{Counting-register distribution, total $1.000000000000$.}\label{tab:shor}
\end{table}

\noindent
Continued fractions on $4/16$ and $12/16$ give $r = 4$; then
$\gcd(7^2-1,15) = 3$ and $\gcd(7^2+1,15) = 5$.

The zeroes in Table~\ref{tab:shor} are the point. Twelve of the sixteen outcomes
are operationally reachable --- some derivation of the term produces each of them
--- and carry amplitude zero. By Definition~\ref{def:deficit} they constitute the
interference deficit $\mathcal{D}(P)$, and by Corollary~\ref{cor:adequacy} the
inclusion $\supp(\sem{P}) \subseteq \Reach(P)$ is strict by exactly those twelve.
Shor's algorithm, in this semantics, \emph{is} a large interference deficit
arranged on purpose.

\subsection{The hazard: unnormalised clauses buy post-selection}\label{ssec:hazard}

\S\ref{ssec:fork} chose the unnormalised reading, with the Born rule applied at
observation. That was the right call for the degenerate race, but nothing in
\S\S\ref{sec:design}--\ref{sec:semantics} requires a clause matrix to be an
isometry. A non-isometric clause shrinks the norm, and renormalising at
observation is precisely post-selection on the surviving branch.

\begin{proposition}[Overshoot]\label{prop:postbqp}
With $\V = \Cx$ and arbitrary clauses, the graded fragment with the Born reading
decides $\mathsf{PostBQP}$, and hence $\mathsf{PP}$.
\end{proposition}

\begin{proof}[Proof sketch]
Lemmas~\ref{lem:1q}--\ref{lem:2q} give universal quantum circuits. The clause
$\mathrm{diag}(1,\epsilon)$ is expressible and non-isometric; applying it and
normalising at observation implements post-selection on the first outcome up to
error $O(\epsilon)$. Conversely the model is simulable with post-selection.
Aaronson's identification of $\mathsf{PostBQP}$ with $\mathsf{PP}$ gives the
second equality.
\end{proof}

Measured, to show that the semantics charges nothing for it: a Hadamard layer on
$n$ qubits followed by $\mathrm{diag}(1,10^{-3})$ on each wire leaves a surviving
norm of $2.5\times 10^{-1}$, $1.25\times 10^{-1}$, $6.25\times 10^{-2}$,
$3.13\times 10^{-2}$ for $n = 4,6,8,10$ --- falling as $2^{-n/2}$ --- while the
renormalised probability of the all-zero string is
$0.999996$, $0.999994$, $0.999992$, $0.999990$. The filter is free.

So the fragment as specified in \S\ref{sec:semantics} is not a quantum
programming language. It is a post-selected one, and it computes $\mathsf{PP}$
--- far more than Shor, and far more than any device. The containment condition
must be stated.

\begin{definition}[Isometric clause]\label{def:isometric}
Let $r$ be a receipt site, let $I$ index the input configurations on which $r$ is
enabled, and let $O$ index the outcomes. The clause $\psi$ at $r$ is
\emph{isometric} when the matrix $A_{io} := \sum_{\sigma \in \Cand:\,
\Out(\sigma)=o} \sem{\psi}(\sigma)$, taken at input $i$, satisfies
$A^{\dagger}A = \mathbb{1}$.
\end{definition}

\begin{proposition}[Decidability of the condition]\label{prop:decidable}
For a clause of the normal form of Theorem~\ref{thm:dnf} with constant
$v_i \in V$ and crisp $\chi_i$, the matrix $A$ is a finite table computable from
the clause and the site's arity, and isometry is decidable by inspection. For
clauses using the graded modality $\langle K\rangle_\V$ the entries depend on
the reachable successors and no such decision procedure is immediate.
\end{proposition}

This gives a natural stratification, and it is the same one extraction wants:
\emph{the modality-free fragment with isometric clauses is the extractable
fragment}, and it is the fragment a compiler to circuits can consume. Clauses
using $\langle K\rangle_\V$ remain available for the stochastic and
inference-driven regimes of \S\ref{sec:pln}, where no unitarity is expected of
them.

\subsection{What the example shows, and what it does not}

It shows that expressiveness is not the obstacle: the fragment is universal for
quantum circuits, and the classical portion of a quantum algorithm is written as
ordinary rholang.

It does not show that graded rho has a quantum idiom of its own. Every gate above
is a hand-built join gadget, and mobility, reflection and name-passing sit idle
throughout; the encoding is circuit-shaped, and a circuit-shaped encoding
demonstrates capability rather than fit. The one primitive the design offers that
a circuit does not hand over directly is Theorem~\ref{thm:scattering}: an
$m$-fold join is a BosonSampling instance in a single term former. Since
candidate sets are permutation sets, the ambient symmetry group of the design is
$S_m$ rather than $\mathbb{Z}_N$, which suggests that a native algorithm --- if
there is one --- would sit nearer the non-abelian hidden subgroup problem than
near Shor. We record the direction without claiming it.

Finally, writing Shor is not running Shor. Evaluating the denotation classically
is exponential; the simulator used for this section is a state-vector simulator
wearing a process calculus. The value of the term is as a \emph{compilation
source}, and Proposition~\ref{prop:decidable} says exactly which terms a
compiler can accept.

%%%%%%%%%%%%%%%%%%%%%%%%%%%%%%%%%%%%%%%%%%%%%%%%%%%%%%%%%%%%%%%%%%%%%%%%%%%%%%%
\section{PLN-valued where-clauses}\label{sec:pln}

\subsection{Why this slot invites PLN}

MeTTaIL describes finitely presentable graph-structured lambda theories, and the
control-flow layer is a rholang whose channels are typed by those language
definitions. Probabilistic Logic Networks are the uncertain-inference framework
developed for the OpenCog and Hyperon line of work, in which knowledge is
annotated with truth values and inference rules carry truth-value formulas
computing a conclusion's annotation from its premises'.

The proposal of this note opens a door between them. If a \texttt{where} clause
is valued in an algebra rather than in the Booleans, then a clause may be valued
\emph{by an inference}: the propensity with which a communication fires is the
truth value PLN currently assigns to the proposition that it should. Execution
stops being gated by belief and starts being \emph{driven} by it. The scheduler
and the reasoner become the same loop.

This is a different use of the machinery from the quantum one, and it is worth
being explicit that it is a different \emph{regime} rather than a different
mechanism.

\subsection{PLN truth values, briefly}

We use PLN's \emph{simple truth values}: a pair $(\str,\conf) \in [0,1]^2$ of a
strength, the estimated probability that a statement holds, and a confidence,
reflecting how much evidence stands behind that estimate.\footnote{Confidence is
formalised in several ways across the literature, including as the width of a
confidence interval around the mean of a second-order distribution; the
NARS-style count form $\conf = n/(n+k)$ for an evidence count $n$ and a personality
parameter $k$ is the one easiest to read operationally, and we use it
informally.} The independence-based deduction formula for term-logic transitivity
is
\[
  \str_{AC} \;=\; \str_{AB}\str_{BC}
  \;+\; \frac{(1-\str_{AB})(\str_C - \str_B \str_{BC})}{1 - \str_B},
\]
with a concept-geometry variant replacing the independence assumption by an
assumption about the shape of concepts,
$\str_{AC} = \str_{AB}\str_{BC} / \min(1, \str_{AB}+\str_{BC})$. Induction and
abduction are obtained by chaining deduction with Bayes' rule. The quantity
$\pow := \str \times \conf$, sometimes called the power of an uncertain
statement, behaves multiplicatively under deduction in the regime where term
probabilities are highly uncertain: $\pow_{AC} \approx \pow_{AB}\,\pow_{BC}$.

\subsection{The resolution algebra, honestly}

The design of \S\ref{sec:design} requires a commutative semiring. PLN's
operations do not straightforwardly furnish one, and it is better to say where
the difficulty lies than to paper over it.

Multiplication is unproblematic. Conjunction of independent statements
multiplies strengths, and confidences compose multiplicatively under the usual
independence assumption, so $\tensorv$ may be taken componentwise; under the
high-uncertainty regime this agrees with deduction on the power.

Addition is the problem. PLN's operation for combining two estimates of the same
proposition is \emph{revision}, a confidence-weighted average
\[
  \str \;=\; \frac{\str_1\conf_1 + \str_2\conf_2}{\conf_1+\conf_2},
\]
with an accompanying rule for the revised confidence. Revision is not the
addition of a semiring: it is idempotent on equal arguments, it does not have
$(0,\cdot)$ as a two-sided unit in the required sense, and it does not distribute
over componentwise multiplication. Taking $\sumv$ to be revision therefore makes
the one-step operator \eqref{eq:onestep} sensitive to the order in which
candidates are aggregated, which is precisely the order the calculus refuses to
fix.

Two responses are available, and we recommend the first.

\begin{enumerate}[leftmargin=1.4em,itemsep=3pt]
\item \textbf{PLN supplies values; $\Rnn$ supplies the algebra.} Let the graded
sort contain ground formulae whose evaluation is a PLN inference, and let the
clause's value be the resulting power $\pow = \str\times\conf \in [0,1]$, with
$\sumv = +$ and $\tensorv = \times$ over $\Rnn$. Then $\V = \Rnn$, the semiring
laws hold, and the whole apparatus of \S\S\ref{sec:design}--\ref{sec:semantics}
applies unchanged. Deduction's multiplicativity on the power means that chaining
inference and chaining communication agree, which is the coherence one wants.
The confidence coordinate is not discarded: it enters the propensity, so a
poorly-evidenced belief yields a rarely-taken branch rather than a forbidden one.
\item \textbf{Carry $(\str,\conf)$ and weaken the algebra.} Keep the pair as the
carrier and accept a lax structure in which aggregation is revision. This is more
faithful to PLN but requires a canonical aggregation order, which must then be
justified as not observable --- a substantial obligation, and one that
reintroduces exactly the kind of implementation-visible choice the design was
meant to abstract.
\end{enumerate}

\begin{principle}[Fuzzy inference drives execution]\label{prin:pln}
Under the first response, a MeTTaIL rewrite fires with propensity equal to the
power PLN currently assigns to the proposition, expressed as a graded ground
formula in the rule's \texttt{where} clause, that the rewrite is warranted here.
Because clauses are re-evaluated at every candidate against the current term
(Proposition~\ref{prop:plastic}), revisions of belief take effect on the next
step without any separate update mechanism.
\end{principle}

\subsection{Three regimes, one syntax}

The point of Table~\ref{tab:algebras} is now visible. The same slot, with the
same formula language and the same per-match evaluation, gives:

\begin{itemize}[leftmargin=1.4em,itemsep=2pt]
\item $\Bool$ --- rholang 1.4, unchanged, by Theorem~\ref{thm:conservativity};
\item $\Rnn$ --- stochastic execution, with PLN as a distinguished source of
values, in which inference schedules computation;
\item $\Cx$ --- coherent execution with interference, in which the clause is a
scattering matrix by Theorem~\ref{thm:scattering}.
\end{itemize}

It is worth stating the boundary between the second and third plainly, since it
is easy to blur. PLN values live in $[0,1]$ and carry no phase; PLN-driven
execution is therefore stochastic and exhibits no cancellation. The interference
regime requires additive inverses in the value algebra, which is a property of
$\Cx$ and not of any probability-valued logic. What the two share is the
mechanism; what separates them is whether $\sumv$ can cancel.

\begin{remark}[A speculation worth naming, and not indulging]\label{rem:paraconsistent}
PLN has been given paraconsistent formulations, and paraconsistent logics are
one of the places complex or otherwise non-order-theoretic truth values have
been contemplated. Whether a paraconsistent PLN would furnish a value algebra
with cancellation --- and hence let inference and interference share not merely a
slot but a semantics --- is a real question and not one this note attempts.
\end{remark}

%%%%%%%%%%%%%%%%%%%%%%%%%%%%%%%%%%%%%%%%%%%%%%%%%%%%%%%%%%%%%%%%%%%%%%%%%%%%%%%
\section{A second worked example: a mortal forager}\label{sec:forager}

\S\ref{sec:shor} exercised the coherent regime. This section exercises the
inference-driven one, and it is drawn from a setting where graded clauses have
somewhere to go rather than being demonstrated in the abstract: the mortal
scientist programme, in which a learner's hypotheses are already installed in
\texttt{where} guards and every observation is charged against a finite token
stack. There, a hypothesis in operational position is an experiment. Here we
value that guard in PLN and ask what changes.

The example is deliberately small --- one forager, one predicate, one belief ---
and every number below is measured by a simulator that runs the clause as
specified.

\subsection{The setting}

A forager meets a stream of specimens. Each is \emph{rich}, carrying a stack of
$\sigma_R = 120$ tokens, or a \emph{decoy}, carrying $\sigma_D = 0$. Breaking one
open costs $\kappa_{\rm break} = 40$, so opening a decoy is a pure loss. A
structural predicate $\chi$ at depth two costs $\kappa_{\rm sense} = 2^2 = 4$ to
evaluate and is \emph{evidence, not a decision procedure}: it is satisfied by
$80\%$ of rich specimens and $30\%$ of decoys. The forager may instead graze a
wall yielding $5$, which is the outside option. It starts with $120$ tokens and
meets at most $400$ specimens; death is failure to fund.

Two worlds. In world $A$ the base rate of rich specimens is $0.40$, so
$\Pr[\text{rich} \mid \chi] = 0.64$ and opening a $\chi$-specimen is worth
$+36.80$ in expectation. In world $B$ the base rate is $0.15$ and $\chi$ is
satisfied by $40\%$ of decoys, so $\Pr[\text{rich}\mid\chi] = 0.26$ and opening
is worth $-8.70$. The forager is not told which world it is in.

\begin{lstlisting}
new bel, prey, wall in {
  bel!( (1, 0.5) ) |                    // (n, s): one unit of prior evidence

  // the two branches race for the belief and the specimen
  for( @(n,s) <- bel ; @spec <- prey  where  Chi(spec) * Pow(n,s) ) {
    Break!(spec) | bel!( revise(n, s, spec) )
  } |
  for( @(n,s) <- bel ; @spec <- prey  where  g ) {
    Graze!(*wall) | bel!( (n, s) )
  }
}
\end{lstlisting}

\noindent
\texttt{Chi} is a crisp structural predicate entering through the indicator
embedding $\crisp{-}$; \texttt{Pow(n,s)} is a graded ground formula evaluating to
the PLN power $\str \cdot \conf$ with $\conf = n/(n+k)$; \texttt{*} is the graded
conjunction $\tensorv$; and \texttt{g} is a scalar constant of $\V = \Rnn$. By
\S\ref{ssec:permatch} both clauses are evaluated per candidate, so the forager's
propensity to open \emph{this} specimen is $\str\conf$ against the wall's $g$.

\begin{remark}[The belief is a message]\label{rem:belief-message}
The belief sits on a channel as data. It is therefore in the tuple space like
any other datum, which means it can be read, raced for, and taken. In the
ecological reading this is not an implementation detail: testimony and the theft
of a hypothesis are the same operation as predation, differing only in which
namespace the name lives in.
\end{remark}

\subsection{What the graded clause replaces}

The assay in the mortal scientist programme is a \emph{complementary guard pair}:
one receipt guarded by $\varphi$ reporting confirmation, a second guarded by
$\neg\varphi$ reporting refutation, and a third outcome $\bot$ when neither fires
before the budget is exhausted. The pair is needed because a Boolean guard that
does not fire is indistinguishable from one still waiting.

This has a dependency worth naming: the second guard requires $\neg$ \emph{in the
hypothesis language}, and a generated logic supplies negation only when the
target logic specification does --- which several members of the generating
family do not.

\begin{proposition}[The trichotomy is the corners of the truth-value square]
\label{prop:trichotomy}
Valuing the guard at a PLN simple truth value $(\str,\conf)$ replaces the pair by
one clause and its complement in the \emph{value algebra}: the confirming branch
carries $\str\conf$, the refuting branch $(1-\str)\conf$, and neither fires with
weight $1 - \conf$. Confirmation, refutation and $\bot$ are then the corners
$\conf = 1, \str = 1$; $\conf = 1, \str = 0$; and $\conf = 0$ of the unit square,
and $\bot$ becomes a coordinate rather than a residue.
\end{proposition}

The negation the assay needs has moved from the formula language, where it may
not exist, to $[0,1]$, where $1 - \str$ always does. This is the same observation
as \S\ref{sec:design}'s --- that negation does not survive to the graded sort ---
read in the direction where it is a gain rather than a cost.

\begin{remark}[Confidence needs an outside option to be visible]
\label{rem:cancel}
If the competing branch is the complement of the same belief, the confidence
coordinate is a common factor: $\str\conf / (\str\conf + (1-\str)\conf) = \str$
exactly, and $\conf$ cancels. Confidence changes behaviour only against an
alternative whose weight does not scale with it. In this example that alternative
is the wall, and in the ecology generally it is autotrophy. Grading a guard buys
nothing unless the learner has something else to do.
\end{remark}

\subsection{The bootstrap trap}

\begin{proposition}[Zero is a fixed point of the power]\label{prop:bootstrap}
A belief with no evidence has $\conf = 0$, hence power $0$, hence a clause value
of $0$, hence --- by the support reading of Remark~\ref{rem:oneslot} --- a
receipt that never fires. It therefore never acquires evidence.
\end{proposition}

Measured: a forager initialised with $n = 0$ opens $0.00$ specimens over $400$
encounters in either world and ends with $520.0$ tokens, exactly the grazing
baseline, against $3543.7$ for the same forager given a single unit of prior
evidence in world $A$. The trap is not death but stasis.

The remedy is available inside the clause language and is worth stating as such:
either the belief carries prior evidence $n_0 \geq 1$, which is PLN's own
convention, or the clause is $\mathrm{Pow}(n,s) \sumv \varepsilon$ for a scalar
$\varepsilon$ --- \emph{exploration is a disjunct}. We use the former below.

\subsection{Two worlds, four foragers}

Four arms. \textsc{naive} has no prior evidence. \textsc{crisp} is the mechanism
as it stands: the guard is $\chi$, weight $1$ when satisfied, confidence
invisible. \textsc{individual} holds a hypothesis about the specimen, so its
evidence count cannot pass $1$ --- the specimen is destroyed by the assay.
\textsc{pln} holds a class hypothesis and accumulates. All at $k = 4$, over
$20{,}000$ lives each.

\begin{table}[h]
\centering
\small
\begin{tabular}{@{}llrrrrr@{}}
\toprule
world & arm & survival & terminal & lifetime & opened & decoys \\
\midrule
$A$ & \textsc{naive} & $1.0000$ & $520.0$ & $400.0$ & $0.00$ & $0.00$ \\
$A$ & \textsc{crisp} & $0.8501$ & $\mathbf{4084.7}$ & $341.6$ & $113.83$ & $40.94$ \\
$A$ & \textsc{individual} & $0.9478$ & $2117.3$ & $381.2$ & $50.68$ & $18.20$ \\
$A$ & \textsc{pln} & $\mathbf{0.9485}$ & $3543.7$ & $381.6$ & $95.53$ & $34.36$ \\
\midrule
$B$ & \textsc{naive} & $1.0000$ & $\mathbf{520.0}$ & $400.0$ & $0.00$ & $0.00$ \\
$B$ & \textsc{crisp} & $0.0029$ & $27.4$ & $30.5$ & $8.72$ & $6.44$ \\
$B$ & \textsc{individual} & $0.2121$ & $95.2$ & $166.8$ & $13.65$ & $10.10$ \\
$B$ & \textsc{pln} & $0.0509$ & $37.2$ & $125.2$ & $14.85$ & $10.99$ \\
\bottomrule
\end{tabular}
\caption{Four foragers, two worlds, $k=4$.}\label{tab:forager}
\end{table}

Two things in this table are worth stating as findings rather than data.

First, in world $A$ the crisp forager has the \emph{highest} mean terminal stack
and the \emph{lowest} survival. The confidence coordinate does not buy yield; it
buys variance. A forager that acts on a satisfied guard from its first encounter
does better on average and dies more often, which is the trade a mortal learner
is actually facing and which a Boolean guard cannot express.

Second, and more sharply: in world $B$ the crisp forager is essentially extinct
at $0.29\%$ survival. It cannot do otherwise. Its policy is the formula, and the
formula does not change, so \emph{a Boolean guard has no coordinate in which to
record disappointment}. The graded forager's policy is the formula's
\emph{value}, which moves. Here is a single life in world $B$:

\begin{table}[h]
\centering
\small
\begin{tabular}{@{}lrrrr@{}}
\toprule
opening & $\str$ & $\conf$ & power & $\Pr[\text{open}]$ \\
\midrule
1 & $0.7500$ & $0.3333$ & $0.2500$ & $0.3333$ \\
2 & $0.5000$ & $0.4286$ & $0.2143$ & $0.3000$ \\
3 & $0.3750$ & $0.5000$ & $0.1875$ & $0.2727$ \\
5 & $0.2500$ & $0.6000$ & $0.1500$ & $0.2308$ \\
\bottomrule
\end{tabular}
\caption{Learning to stop. Strength falls as decoys accumulate; confidence
rises; the product falls, and with it the propensity to open.}\label{tab:trace}
\end{table}

\begin{remark}[Exploration and exploitation are not scheduled]
\label{rem:explore}
Nothing in the clause is a schedule. Early behaviour is cautious because $\conf$
is small, late behaviour is decisive because $\conf$ is large, and the transition
rate is set by how fast evidence arrives. The usual hyperparameter is absent
because the quantity it approximates is a coordinate of the truth value.
\end{remark}

\subsection{Amortisation, measured}

The class hypothesis and the individual hypothesis differ in exactly one place:
the individual's evidence count cannot grow, because the assay destroys its
subject. That is not a difference in what can be believed but in what can be
\emph{accumulated}, and PLN's confidence is the ledger of accumulation.

Measured over survivors in world $A$ at $k=4$, against the true posterior
$0.6400$:

\begin{center}
\small
\begin{tabular}{@{}lrrr@{}}
\toprule
arm & mean $|\str - \text{truth}|$ & mean evidence count & survivors \\
\midrule
\textsc{individual} & $0.2039$ & $2.00$ & $18{,}917$ \\
\textsc{pln} & $\mathbf{0.0383}$ & $101.54$ & $18{,}900$ \\
\bottomrule
\end{tabular}
\end{center}

\noindent
A factor of $5.3$ in accuracy, for evidence that costs nothing extra --- the
openings happen either way; only the bookkeeping differs. This is the
amortisation advantage in its defensible form. It says that reusable class
hypotheses recover their acquisition cost across instances and individual ones
cannot, which is a selective advantage; it does not say that individual
hypotheses have no value, and the table shows they do not (the
\textsc{individual} arm outlives \textsc{crisp} in both worlds).

\subsection{The personality parameter acquires a metabolic optimum}

In PLN and its relatives, $k$ is a free parameter: the amount of evidence at
which confidence reaches one half. It is chosen by the designer's taste. In a
mortal ecology it is not free, because it is priced on both sides. A credulous
forager (small $k$) becomes confident before it has grounds and opens decoys it
cannot afford; a sceptical one (large $k$) never becomes confident enough to eat.

Scored by mean $\log$ terminal wealth --- the criterion appropriate to a
multiplicative, absorbing process, where ruin contributes zero rather than a
small number --- the optimum is interior. World $A$, $40{,}000$ lives per point:

\begin{table}[h]
\centering
\small
\begin{tabular}{@{}lrrrrrrrrrrr@{}}
\toprule
$k$ & $4$ & $5$ & $6$ & $7$ & $8$ & $\mathbf{9}$ & $10$ & $12$ & $14$ & $16$ & $20$ \\
\midrule
survival & $.944$ & $.952$ & $.959$ & $.967$ & $.970$ & $.976$ & $.979$ & $.984$ & $.988$ & $.991$ & $.995$ \\
$\log$-fitness & $7.744$ & $7.785$ & $7.818$ & $7.853$ & $7.852$ & $\mathbf{7.874}$ & $7.872$ & $7.858$ & $7.832$ & $7.798$ & $7.715$ \\
\bottomrule
\end{tabular}
\caption{An optimal degree of scepticism, $k^\ast = 9$.}\label{tab:ksweep}
\end{table}

\noindent
In an ensemble where each life draws world $A$ or world $B$ with equal
probability --- the honest setting, since a learner does not know which world it
is in --- the optimum moves to $k^\ast = 64$, with $\log$-fitness $6.343$,
$\mathbf{6.431}$, $6.354$, $6.316$, $6.263$ at $k = 32, 64, 96, 128, 512$. In
world $B$ alone it runs to the boundary, because when the evidence is a trap the
best policy is abstinence and large $k$ is how the forager approximates it.

\begin{principle}[Scepticism is priced by the environment]\label{prin:k}
$k^\ast$ is not a property of the reasoner. It rises with the risk of the world
the reasoner is in --- here by a factor of about seven between a world whose
evidence is sound and an ensemble containing a trap. A parameter that uncertain
inference leaves to taste is, in a mortal setting, selected.
\end{principle}

We regard this as the example's main return, and note that it runs in both
directions: the graded clause gives the ecology a principled exploration
schedule, and the ecology gives PLN a selection criterion for a parameter it
otherwise has to stipulate.

\subsection{Up the fractal}\label{ssec:fractal}

The forager is a single learner. The programme this example is drawn from
continues upward --- to learners that compose by parallel composition over shared
namespaces, to lineages, and to populations --- and the graded clause has
something to say at each level. We record two observations and claim neither.

The first is that PLN's operation for combining two estimates of the same
proposition is \emph{revision}, which adds evidence counts. Two learners that
share a hypothesis namespace and exchange truth values therefore have
superadditive confidence, in a setting where the same programme's own results
make metabolism subadditive: tokens are conserved and formulae are not. Testimony
would then be the operation at which a composite strictly exceeds the sum of its
parts, and PLN would say by how much.

The second is the hazard attending the first. Revision assumes the evidence being
combined is independent. Two learners who both heard from a third double-count,
and confidence rises with no evidence behind it. \emph{Belief inflation is the
epistemic analogue of minting} --- and where the conservation law of a token
economy forbids minting outright, nothing in the epistemic layer forbids this.
Whether a namespace discipline can restore conservation to the evidence ledger
seems to us the right question to ask next.

Both of these need the truth value carried as a pair rather than collapsed to its
power, which puts them outside the algebra recommended in \S\ref{sec:pln} and
into the lax structure set aside there. They belong to a subsequent paper, along
with the back-patching of the results above into the programme they came from.

%%%%%%%%%%%%%%%%%%%%%%%%%%%%%%%%%%%%%%%%%%%%%%%%%%%%%%%%%%%%%%%%%%%%%%%%%%%%%%%
\section{Implementation}\label{sec:impl}

The design is deliberately small on the implementation side, and most of what it
needs already exists.

\paragraph{Candidate enumeration.} The store presentation gives the redex
enumerator directly: the map from channels to rows of data and waiting
continuations, together with the join index, is precisely $\Cand$. The one
requirement the current representation does not meet is occurrence identity ---
per-channel occupancy counts collapse the distinction between congruent data, and
identifiers allocated during firing are not stable across branches. Candidates
must be addressed by (channel, canonical slot), deterministically and
branch-independently. This is the same addressing obligation that arises in the
compilation route, and it should be solved once.

\paragraph{Basis identity.} Interference is addition of coefficients on a common
basis vector, so the implementation needs a canonical name for a term modulo
structural congruence. The content-addressed store already supplies one. This
also repairs a known defect: a configuration fingerprint that formats values to
fixed decimal precision cannot distinguish $1$, $-1$ and $i$, and therefore
cannot represent the very distinctions the grading exists to make.

\paragraph{Clause evaluation.} The crisp sort is the existing condition language
and its evaluator; the graded sort adds scalars, two binary operations and the
graded modality. The modality is the only construct requiring successor
exploration, and it takes the inference budget.

\paragraph{Two readings out.} A simulator should emit both the support/Born
reading and the coherent sum, since by Corollary~\ref{cor:adequacy} they differ
exactly on the interference deficit. If they never differ on a corpus of test
theories, the grading is doing nothing and the implementation is wrong.

\paragraph{What retires.} Partition construction and its checked funnel; the
requirement that keys be structural, as a precondition for partition locality;
the update-function plumbing and the declarative update sub-language it wanted;
and the class-global phase, which was in any case observationally inert.

%%%%%%%%%%%%%%%%%%%%%%%%%%%%%%%%%%%%%%%%%%%%%%%%%%%%%%%%%%%%%%%%%%%%%%%%%%%%%%%
\section{Open threads}\label{sec:open}

\begin{obligation}[Length-asymmetric reconvergence]\label{ob:asym}
Theorem~\ref{thm:scattering} secures interference at a single join. The general
question is which closed loops survive the causal quotient of
Definition~\ref{def:causal} in pure rho --- that is, whether two causally
inequivalent derivations of different lengths can reach a common term. Reflection
is the natural suspect, since dropping a quoted process makes the residual
communication count depend on which datum won. A negative answer would confine
interference to joins; a positive one would open a second, non-local source.
\end{obligation}

\begin{obligation}[Normalisation]\label{ob:norm}
\S\ref{ssec:fork} recommends the unnormalised reading with Born at observation.
The alternative branches should be worked out far enough to confirm that the
choice is a choice of observable rather than of correctness.
\end{obligation}

\begin{obligation}[Enforcing isometry]\label{ob:unitary}
\emph{Upgraded by \S\ref{ssec:hazard} from an open thread to a containment
requirement.} Definition~\ref{def:isometric} states the condition and
Proposition~\ref{prop:decidable} says it is decidable on the modality-free
fragment. What remains is to turn that into a static discipline --- a judgement
on clauses, checked at elaboration, refusing non-isometric clauses in a
declared coherent mode while leaving them available in the stochastic and
inference-driven modes where no unitarity is expected. Until that exists, a
graded program should be read as post-selected.
\end{obligation}

\begin{obligation}[Adaptivity and universality]\label{ob:klm}
A continuation may install further joins depending on what it received, which is
adaptivity. Whether adaptive graded joins reach a universal fragment --- the
question KLM answers affirmatively for linear optics --- is the question on which
a verdict about advantage ultimately rests.
\end{obligation}

\begin{obligation}[Categorical status]\label{ob:cat}
The construction looks like a Kleisli category for the free-module monad of $\V$
over terms modulo congruence, with the graded modality as the structure map of a
coalgebra. Making this precise would connect the design to the existing
coalgebraic account of the store denotation and its two-layer full abstraction.
\end{obligation}

\begin{obligation}[Testimony, and conservation of evidence]\label{ob:fractal}
\S\ref{ssec:fractal}. Carrying $(\str,\conf)$ rather than the power puts
revision back in play, which makes confidence superadditive under testimony and
opens the double-counting failure --- belief inflation as the analogue of
minting. Both require the lax algebra set aside in \S\ref{sec:pln} and belong to
a subsequent paper, together with back-patching \S\ref{sec:forager}'s results
into the mortal-scientist programme they are drawn from.
\end{obligation}

\begin{obligation}[Paraconsistent values]\label{ob:para}
Remark~\ref{rem:paraconsistent}: whether an inference calculus can supply a value
algebra with cancellation.
\end{obligation}

%%%%%%%%%%%%%%%%%%%%%%%%%%%%%%%%%%%%%%%%%%%%%%%%%%%%%%%%%%%%%%%%%%%%%%%%%%%%%%%
\section{Conclusion}\label{sec:conc}

The rho calculus does not say how its races are settled, and that silence is a
parameter rather than a gap. Filling it differently gives different executions of
the same program, and a programmer who wants to arrange interference needs to be
able to write the filling down.

The change proposed here is one line of grammar: let the \texttt{where} clause of
the for-comprehension take values in a commutative semiring, and evaluate it once
per candidate match rather than once per communication. Boolean values give back
today's language exactly. Non-negative reals give stochastic execution, with
probabilistic-logic inference as a natural source of values, so that what the
system believes decides what it does next. Complex values give interference, and
the clause becomes a scattering matrix whose permanent is the amplitude of the
symmetric outcome.

Shor's algorithm fits, and fitting it was informative in an unexpected
direction. Universality is easy, the classical portion of the algorithm is
ordinary rholang, and the algorithm's characteristic behaviour shows up in this
semantics as a large interference deficit arranged on purpose --- twelve of
sixteen reachable outcomes carrying amplitude zero. What the exercise exposed is
that the fragment as first specified reaches too far: with clause matrices
unconstrained, the renormalisation performed at observation is post-selection,
and the fragment computes $\mathsf{PP}$. The isometry condition is therefore not
a compilation convenience but the boundary between a quantum language and a
magic one, and the fragment on which it can be checked --- clauses without the
graded modality --- is precisely the fragment a circuit compiler can accept.

The inference-driven regime was exercised too, on a mortal forager whose
hypothesis is a PLN truth value sitting in a guard. Three things came out of it.
The complementary guard pair such a learner needs in order to tell refutation
from patience turns out to require negation in the hypothesis language, which a
generated logic need not supply --- and grading moves that negation into $[0,1]$,
where it always exists. A Boolean guard has no coordinate in which to record
disappointment, so it cannot learn to stop, and in a world whose evidence is a
trap it does not survive. And PLN's personality parameter, which uncertain
inference leaves to the designer's taste, acquires an interior optimum that rises
with the risk of the environment: scepticism, priced.

Two things about this were not evident at the outset. The first is that weight
maps were not wrong but merely linear: they are the normal forms
$\bigoplus_i v_i \tensorv \crisp{\chi_i}$, and the partition discipline that
attended them was the price of making a lookup single-valued rather than a fact
about weights. The second is that the recombination interference requires ---
the half of the problem no weight map addressed --- was already in the language.
A join over a contended channel has one candidate per assignment and one outcome
per orbit of the continuation's symmetry, so whether a race interferes is decided
by whether the continuation troubles to tell its inputs apart. That is a dial a
programmer already knows how to turn.

\end{document}
